% Shared exercise chunk: l4-binary-count-accumulation
\begin{exercisechunk}
\item \textbf{Persistent errors accumulate.}
\label[exercise]{ex:binary-count-accumulation}
\exerciseprofile{3}{2}{3}{\exproof\quad\excalculation\quad\exinterpretation}{%
Derive exact and asymptotic error counts and explain the quadratic accumulation
caused by persistent errors.}
Use the setting and Markov generators from \cref{ex:binary-count}.
\begin{enumerate}[(a),beginpenalty=10000]
\item Prove
\begin{align}
 \E_{Q_{\mathrm{ind}}}[c_{\mathrm{cnt}}]&=T\varepsilon,
 \label{eq:ind-cost}\\
 \E_{Q_{\mathrm{bad}}}[c_{\mathrm{cnt}}]
 &=\sum_{t=1}^T[1-(1-\varepsilon)^t]
 =T-\frac{(1-\varepsilon)[1-(1-\varepsilon)^T]}{\varepsilon}.
 \label{eq:bad-cost}
\end{align}
Obtain the second identity both from marginal error probabilities
and by conditioning on the first error position. Also prove
\begin{equation}
 \E_{Q_{\mathrm{ind}}}[c_{\mathrm{fin}}]=\varepsilon,\qquad
 \E_{Q_{\mathrm{bad}}}[c_{\mathrm{fin}}]=1-(1-\varepsilon)^T.
 \label{eq:terminal-cost}
\end{equation}
\item Prove the finite bounds
\[
 \frac{\varepsilon T(T+1)}2
 -\frac{\varepsilon^2T(T+1)(T-1)}6
 \le\E_{Q_{\mathrm{bad}}}[c_{\mathrm{cnt}}]
 \le\frac{\varepsilon T(T+1)}2.
\]
Deduce \(\E_{Q_{\mathrm{bad}}}[c_{\mathrm{cnt}}]
\sim\varepsilon T(T+1)/2\) when \(T\varepsilon\to0\).
Explain why the expected count can grow quadratically with \(T\)
in this regime.
\end{enumerate}
\end{exercisechunk}
